\documentclass[10pt]{article}
\usepackage[margin=0.78in]{geometry}
\usepackage{amsmath,amssymb}
\usepackage[colorlinks=true,linkcolor=blue,citecolor=blue,urlcolor=blue]{hyperref}

\title{A negative answer to the symmetric-matrix TTO question}
\author{Literature-status note for arXiv:2306.12183}
\date{July 21, 2026}

\begin{document}
\maketitle

\noindent\textbf{Status.} Literature already answered.

\section{Source question}

O'Loughlin and Virtanen ask in Question~4.1 on page~12 of [1]:

\begin{quote}
Is every symmetric matrix unitarily equivalent to a direct sum of TTOs?
\end{quote}

Here ``symmetric'' means complex symmetric as a matrix, and TTO abbreviates
truncated Toeplitz operator.  The question is motivated by the possibility of
using TTOs as model operators for numerical ranges and Crouzeix's conjecture.

\section{Later negative answer}

O'Loughlin's July 2026 paper [2] answers this question negatively.
Corollary~4.4 on page~7 states that for every $n\ge10$, the question is false.
More explicitly, for each such $n$ there exists an irreducible complex
symmetric $n\times n$ matrix that is not unitarily equivalent to a truncated
Toeplitz operator.

The proof uses semialgebraic dimension.  Let $\mathcal{CS}_n$ be the real
vector space of complex symmetric $n\times n$ matrices and let $\mathcal I$
be its irreducible locus.  The answering paper identifies a semialgebraic
family $\mathcal C$ containing symmetric matrix representations of
$n$-dimensional TTOs and proves
\[
 \dim_{\mathbb R}\mathcal U(\mathcal C)
 \le 7n-6+\frac{n(n-1)}2
 < n(n+1)
 =\dim_{\mathbb R}(\mathcal I\cap\mathcal{CS}_n)
 \qquad(n\ge10).
\]
Thus not every irreducible symmetric matrix has a TTO model.  An irreducible
matrix cannot be unitarily equivalent to a nontrivial direct sum, so this also
rules out a direct sum of TTOs and answers the source question.

\section{Scope}

The result is existential: it proves that counterexamples form a
dimension-theoretically unavoidable family but does not display a particular
matrix.  It settles the universal question negatively; the individual sizes
$4\le n\le9$ remain open.  The negative answer removes one proposed TTO-model
route but does not disprove or prove Crouzeix's conjecture.

\section*{References}
\small\raggedright
\noindent[1] R.~O'Loughlin and J.~Virtanen,
\emph{Crouzeix's conjecture for classes of matrices},
arXiv:2306.12183v2, 2023.\par\smallskip
\noindent[2] R.~O'Loughlin,
\emph{Semialgebraic Dimension and Truncated Toeplitz Models for Complex
Symmetric Matrices}, arXiv:2607.14019v1, 2026.

\end{document}
